Pour augmenter la réflexion d'une plaque de verre (indice \num{1.52}), on la recouvre d'une mince couche transparente d'indice de réfraction \num{\v{nc}}. On veut que la lumière de longueur d'onde \qty{\v{lam}}{nm} (dans l'air), arrivant sous incidence normale, soit réfléchie le plus fortement possible.

\begin{enumerate}[label=\alph*)]
    \item Déterminez l'épaisseur minimale de la couche.
    \item Déterminez l'épaisseur suivante qui produirait le même effet.
    \item Avec l'épaisseur trouvée en a), quelle longueur d'onde est complètement absente de la lumière réfléchie ?
\end{enumerate}
